\documentclass[../../main.tex]{subfiles}
\begin{document}
\section{奇异上同调}\label{section:ch18-singular-cohomology}

除奇异同调群外，对任意拓扑空间 $M$ 与任意阿贝尔群 $G$，还可以定义一列密切相关的群 $H^p(M;G)$，称为\emph{系数在 $G$ 中的奇异上同调群}。精确的定义对我们的目的并不重要；我们只关心特殊情形 $G=\mathbb{R}$，此时可以证明 $H^p(M;\mathbb{R})$ 是一个实向量空间，并且自然同构于从 $H_p(M)$ 到 $\mathbb{R}$ 的群同态空间 $\operatorname{Hom}(H_p(M),\mathbb{R})$。（为简单起见，我们把这取作 $H^p(M;\mathbb{R})$ 的定义。）任意连续映射 $F:M\to N$ 诱导一个线性映射 $F^*:H^p(N;\mathbb{R})\to H^p(M;\mathbb{R})$。

代数拓扑中有一个重要的定理称为\emph{泛系数定理}，它说明系数在任意群中的奇异上同调群如何从奇异同调群恢复。因此上同调群并不含有同调群尚未编码的新信息；但它以一种不同的方式组织这些信息，在许多场合更为方便。特别地，奇异上同调群与 de Rham 上同调群一样定义反变函子，这使两者更容易比较。

\begin{proposition}[奇异上同调的性质]\label{proposition:ch18-singular-cohomology-properties}
\begin{enumerate}[(a)]
\item 对任意单点空间 $\{q\}$，除 $p=0$ 外 $H^p(\{q\};\mathbb{R})$ 平凡，而 $H^0(\{q\};\mathbb{R})$ 是一维的。
\item 若 $M=\coprod_j M_j$，则 $H^p(M;\mathbb{R})\cong\prod_j H^p(M_j;\mathbb{R})$。
\item 同伦等价的空间有同构的奇异上同调群。
\end{enumerate}
\end{proposition}

我们需要的关于奇异上同调群的关键事实是：它也满足一个 Mayer--Vietoris 定理。

\begin{theorem}[奇异上同调的 Mayer--Vietoris]\label{theorem:ch18-mayer-vietoris-cohomology}
设 $M,U,V$ 满足\cref{theorem:ch18-mayer-vietoris-singular} 的假设。下面的序列正合：
\begin{align*}
\cdots\to H^p(M;\mathbb{R})\xrightarrow{k^*\oplus l^*}H^p(U;\mathbb{R})\oplus H^p(V;\mathbb{R})
\xrightarrow{i^*-j^*}H^p(U\cap V;\mathbb{R})\xrightarrow{\delta^*}H^{p+1}(M;\mathbb{R})\to\cdots.
\end{align*}
\end{theorem}
\end{document}
